% Transcription — fonds Alexandre Grothendieck, shelfmark 60, batch 4 (pages 61-69).
% Established from the facsimile archives/batches/60/batch-04.pdf (PDF page k =
% archivists' page 60+k, no cover sheet), cut from a copy of 60.pdf fetched
% through the project's relay
% (https://grothendieck-relay.onrender.com/source/60.pdf); tiles in
% archives/tiles/60/p61-p69. Per the skill transcribe-grothendieck. The folder
% is sixty-nine pages in four batches (1-20, 21-40, 41-60, 61-69); this is the
% last, nine pages. Batch 3 (transcripts/60/batch-03.fr.tex) was read first
% and its notation is kept.
%
% Pass: Opus 5.5 (claude-opus-5-5), 2026-09-26 — first pass, unchecked against the pages by a human.
%
% Pages skipped: 61, 63, 65, 67, 69 — the typed versos of the sheets
% (Bourbaki rédaction n° 392, pages 6-45, 6-35, 6-34, 6-33, 6-32: exercises
% of commutative algebra), with nothing in his hand: skipped without
% description by the settled rule.
%
% Pages summarised: none. Pages transcribed: 62, 64, 66, 68 — four rectos in
% ink. His own pagination: none visible (the pencilled numbers bottom left
% are the archivists'). No numbered formulas or labelled runs. P. 64 is
% largely cancelled by three long diagonal strokes and several struck lines;
% it is transcribed and the cancellation recorded once.
%
% What the batch is about. It does not continue the « cohomologie à
% l'infini » of p. 60 directly: it takes up the functional equation in the
% abstract form of p. 54, L(1/qt) = A(t) L(t) for a rational function L,
% A(t) of lower « niveau », and asks whether it can be simplified by a
% change L' = lambda L so that L'(1/qt) = c L'(t), i.e.
% lambda(1/qt)/lambda(t) = c A(t)^{-1} (p. 62); two solutions differ by alpha
% with alpha(1/qt) = a alpha(t), a = ±1 (circled); the group Z/2Z acting by
% t -> 1/qt, the question becomes a cocycle-type equation lambda^g lambda^{-1}
% = ... (top of p. 64). There is a trivial solution lambda = 1/L, excluded
% by requiring lambda of « niveau » strictly below L (p. 64). In a cancelled
% passage he assumes L(0) = 1 and zeros/poles of absolute value p^{-i/2},
% rho <= i/2 <= rho+n, and writes L(t) = L_rho(t) ... L_{rho+n}(t)
% L'_{rho+n-1}(pt) ... L'_rho(p^n t) with the weights of the inverse
% eigenvalues. P. 66: correction by lambdas of the same shape, of weight
% between rho and rho+n but of niveau <= rho+n-1; assuming A(t) of the same
% type and functional equations for each L_i, L'_i, he writes
% L(1/(p^{n+rho} t)) as delta t^chi times a product. P. 68: the quotient
% L(1/qt)/L(t) and A(t) = delta t^chi prod_i L'_i(t)/L'_i(p^{rho+n-i}t)
% prod_i L_i(p^{rho+n-i}t)/L_i(t); the case n = 1 begun and left.
%
% Notation, as in batch 3: q, p, rho; L', lambda, alpha as on the page;
% delta, chi for the constant and exponent of a functional equation (with
% primes and indices as he writes them). Here q = p^{n+rho} is implicit on
% p. 66 (he substitutes 1/(p^{n+rho}t) where p. 68 writes 1/qt); the
% exponent of q on p. 62 is blotted. « niveau » is his word, in quotes.
%
% Readings to check first: the last lines of p. 62 and the top line of p. 64
% (the group action and the equation lambda^g lambda^{-1} = ...); the
% marginal note of p. 64; the arguments p^2 t, p^n t on pp. 64-66.

\documentclass[11pt,a4paper]{article}
% Le préambule commun à toutes les transcriptions du fonds Grothendieck.
%
% Il tient en une idée : **l'incertitude se compose, elle ne se résout pas.**
% Une transcription qui lisse un mot douteux a détruit la seule chose pour
% laquelle on la faisait — un lecteur doit pouvoir savoir, à chaque endroit,
% ce qui était sur le feuillet et ce qui a été ajouté. Les six macros qui
% suivent sont donc l'appareil critique, et non de la décoration.
%
% Les mêmes commandes sont reconnues par `scripts/render.mjs`, qui produit la
% vue de lecture du site. Une macro ajoutée ici doit l'être aussi là-bas,
% faute de quoi le rendu échouera bruyamment — ce qui est voulu : un rendu
% silencieusement faux est pire qu'un rendu absent.

\ProvidesPackage{grothendieck}[2026/08/08 Transcriptions du fonds Grothendieck]

\usepackage{amsmath,amssymb,amsthm}
\usepackage{mathrsfs}
\usepackage{stmaryrd}
% Les diagrammes commutatifs sont partout dans ce fonds : c'est la langue
% même de la géométrie algébrique de Grothendieck, et une transcription qui
% les rendrait en prose ne transcrirait rien.
\usepackage{tikz-cd}
% Français par défaut — c'est la langue des feuillets — mais l'anglais est
% chargé aussi : la traduction et le résumé sont des documents anglais, et la
% typographie française (espaces avant les deux-points) y serait une faute.
% Ces documents basculent par \selectlanguage{english} en tête de corps.
\usepackage[english,french]{babel}
\usepackage{fontspec}
\usepackage{geometry}
\geometry{a4paper,margin=25mm,marginparwidth=28mm}
\usepackage{marginnote}
\usepackage{xcolor}
% ulem, not soul. `soul`'s \st and \ul are built on a character-by-character
% scan that XeLaTeX's font machinery breaks: the document fails with an
% undefined control sequence rather than degrading. ulem does the same two jobs
% and is Unicode-engine safe. `normalem` stops it hijacking \emph, which the
% transcriptions use for Grothendieck's own emphasis.
\usepackage[normalem]{ulem}
\usepackage{hyperref}
\hypersetup{colorlinks=true,linkcolor=black,urlcolor=black,pdfborder={0 0 0}}

% --- Métadonnées du lot -----------------------------------------------------
% Reprises telles quelles par le script de rendu, qui les affiche en tête de
% la vue de lecture. Elles fixent ce que le fichier prétend couvrir : sans
% elles, une transcription flotte sans point d'attache dans un fonds de seize
% mille feuillets.

\newcommand{\folder}[1]{\gdef\grothFolder{#1}}
\newcommand{\batch}[1]{\gdef\grothBatch{#1}}
\newcommand{\leaves}[2]{\gdef\grothFirst{#1}\gdef\grothLast{#2}}
\newcommand{\foldertitle}[1]{\gdef\grothFoldertitle{#1}}
\gdef\grothFolder{?}\gdef\grothBatch{?}\gdef\grothFirst{?}\gdef\grothLast{?}\gdef\grothFoldertitle{}

% --- Le résumé --------------------------------------------------------------
%
% Il ouvre la lecture modernisée, sous le titre « Résumé », et il est composé
% en petit. Ce n'est pas un ornement : c'est le seul passage du document qui
% ne soit pas des mathématiques, et un lecteur doit pouvoir le reconnaître
% avant de l'avoir lu — puis le sauter s'il connaît déjà le sujet, ou s'y
% arrêter s'il ne le connaît pas. Le filet qui le suit marque la reprise du
% corps.

\newenvironment{resume}
  {\small\setlength{\parskip}{0.45em}}
  {\par\medskip\hrule\medskip}

% --- Mots-clés ---------------------------------------------------------------
%
% En anglais, en clôture du résumé de la lecture modernisée : le vocabulaire
% moderne sous lequel chercher ce que les feuillets construisent. Le manifeste
% du site (`npm run manifest`) les extrait pour étiqueter la cote entière —
% cette ligne est la source unique des tags, il n'y a pas de fichier à côté.
\newcommand{\keywords}[1]{\par\smallskip\noindent
  {\footnotesize\textsc{Keywords} --- \textit{#1}}\par}

% --- L'appareil critique ----------------------------------------------------

% Le numéro de feuillet, en marge. C'est l'ancre qui permet de régler un
% désaccord sur une formule en regardant *un* feuillet plutôt qu'une chemise
% de six cents. Le site s'en sert aussi pour tourner les pages du fac-similé
% quand on fait défiler la transcription.
\newcommand{\leaf}[1]{%
  \marginnote{\footnotesize\color{gray!70}\textsf{#1}}%
  \label{leaf:#1}\ignorespaces}

% Illisible : on ne devine pas. Un mot inventé qui se lit comme les autres est
% le pire résultat possible.
\newcommand{\ill}{\textcolor{red!60!black}{[\,\ldots\,]}}

% Lecture douteuse : le mot est donné, et signalé comme incertain.
\newcommand{\uncertain}[1]{\uline{#1}}

% Ajout de l'éditeur — une lettre manquante, un mot sous-entendu.
\newcommand{\add}[1]{\textcolor{blue!50!black}{[#1]}}

% Biffé par Grothendieck lui-même. Ce qu'il a rayé fait partie du document :
% une rature dit souvent où la pensée a changé de direction.
\newcommand{\struck}[1]{\sout{#1}}

% Note du transcripteur, sur l'état matériel du feuillet.
\newcommand{\note}[1]{\par\smallskip{\small\color{gray!60!black}\textit{[#1]}}\par\smallskip}

% Note marginale de Grothendieck, distincte du corps du texte.
\newcommand{\marginal}[1]{\par\smallskip\noindent\hspace{1em}%
  {\small\color{gray!60!black}#1}\par\smallskip}

% --- Titre ------------------------------------------------------------------

\AtBeginDocument{%
  \begin{center}
    {\large\bfseries Fonds Alexandre Grothendieck --- cote n\textsuperscript{o}~\grothFolder}\\[2pt]
    {\itshape\grothFoldertitle}\\[4pt]
    {\small feuillets \grothFirst--\grothLast\ (lot \grothBatch)}\\[2pt]
    {\footnotesize\color{gray!60!black}Transcription établie d'après le fac-similé numérisé
     par l'Université de Montpellier.\\
     Les lectures incertaines sont soulignées, les illisibles notées [\,\ldots\,],
     les ajouts entre crochets.}
  \end{center}
  \vspace{1em}}

\endinput


\folder{60}
\batch{4}
\pages{61}{69}
\dating{[à partir de 1962- à partir de 1971]}
\watermark{Édition de démonstration}
\foldertitle{Fonctions L / Équation fonctionnelle des fonctions L (cas géométrique) : notes manuscrites (s.d.)}

\begin{document}

\page{62}
$L(t)$ une fonction rat.\ en $t$, \struck{\ill{}} $q = p$ \ill{}
\note{l'exposant de $p$ est noyé sous une tache d'encre ; le mot biffé
avant « $q = p$ » commence par un $q$}
\[
  L\Bigl(\frac{1}{qt}\Bigr) = A(t) L(t)
\]
$A(t)$ de « niveau » inférieur.

On veut simplifier l'équation fonctionnelle, en posant
\[
  L'(t) = \lambda(t) L(t)
\]
d'où
\[
  L'\Bigl(\frac{1}{qt}\Bigr) = \lambda\Bigl(\frac{1}{qt}\Bigr) L\Bigl(\frac{1}{qt}\Bigr)
  = \lambda\Bigl(\frac{1}{qt}\Bigr) A(t) L(t)
  = \lambda\Bigl(\frac{1}{qt}\Bigr) \lambda(t)^{-1} A(t) L'(t)
\]
et on aimerait obtenir $L'(\frac{1}{qt}) = c L'(t)$, i.e.
\[
  \frac{\lambda(\frac{1}{qt})}{\lambda(t)} = c A(t)^{-1}
\]
Si on trouve un tel $\lambda$, tout autre $\lambda'(t) = \alpha(t)\lambda(t)$
(\uncertain{est} \uncertain{déterminé} \uncertain{par} \ill{}
\uncertain{solution} \uncertain{générale}
\note{« avec une cte $c'$ » est écrit au-dessus de cette parenthèse, avec
un trait qui la relie à $\lambda'$ ; la parenthèse n'est pas fermée}
\[
  \frac{\alpha(\frac{1}{qt})}{\alpha(t)} = \frac{c'}{c}
  \qquad \text{i.e.} \qquad
  \alpha\Bigl(\frac{1}{qt}\Bigr) = a\,\alpha(t)
\]
avec $a$ une cte, (nécessairement $\pm 1$)
\note{« nécessairement $\pm 1$ » est entouré}
donc dans l'équation fonctionnelle \ill{} \ill{} \ill{} \ill{}
\ill{} \uncertain{mult.}
\note{la fin de cette phrase, à droite, est écrite sur deux lignes,
surchargée et soulignée de plusieurs traits ; une partie semble biffée ;
elle ne se lit pas}

\note{un long trait horizontal sépare ce qui précède de ce qui suit}
Faisons opérer le groupe \struck{\ill{}} $\mathbb{Z}/2\mathbb{Z}$ sur les
fonctions rationnelles $\neq 0$, via la substitution
$t \mapsto \frac{1}{qt}$. Alors \ill{} \ill{} \ill{}) la question
\uncertain{en} \ill{} $\lambda$
\note{un trait vertical, en marge gauche, accompagne les deux dernières
lignes ; la phrase se poursuit en tête de la page 64}

\page{64}
\uncertain{équivaut} à \uncertain{celle} \ill{} \ill{} \ill{} \ill{}
\uncertain{un} \ill{}.
\note{un long trait horizontal passe sur toute cette ligne, qui pourrait
être biffée ; un mot au milieu est biffé à part. Un encadré, relié à la
ligne par un trait, porte l'insertion :}
\[
  \boxed{\underline{\lambda}^{g}\, \lambda^{-1} =}
\]
Il y a une solution triviale
\[
  \lambda = \frac{1}{L} \qquad \text{i.e.} \qquad L'(t) = 1 .
\]
Mais nous devons \struck{que} imposer que $\lambda$ \struck{\ill{}} soit de
« niveau » \uncertain{strict.}\ \uncertain{inf.}\ à $L$.

Pour préciser ce point, nous supposons,
\note{tout le passage qui suit, jusqu'à « de niveau $\leq \rho+n+1$ », est
barré de trois longs traits obliques, et plusieurs de ses lignes sont en
outre biffées d'un trait horizontal ; il est donné tel qu'il se lit}
$L(0) = 1$, que les racines et pôles de $L$ sont des valeurs absolues de
la forme
$p^{-\frac{i}{2}}$ avec $\rho \leq i \leq \rho + 2n$ i.e.\
$\rho \leq \frac{i}{2} \leq \rho + n$,
\note{les deux $\rho$ de la double inégalité en $i$ surchargent un $2\rho$
biffé ; « $L(0) = 1$, » est encadré}
\struck{et \ill{} \uncertain{coefficients} \ill{} \ill{}
\uncertain{puissance} de $p^{\frac{1}{2}}$, dans des \ill{}}
\struck{\uncertain{propre} \ill{} \ill{}}
(et \uncertain{on} \uncertain{écrit} de \uncertain{restreindre}
aux $\lambda$ ayant la \emph{même propriété}, \struck{\ill{} \ill{}
\ill{}} (\uncertain{de} \uncertain{poids} \uncertain{compris} entre
$\rho$ et $\rho+n$) et de plus de « niveau » $\leq \rho+n+1$
\note{« $\rho+n+1$ » est écrit sous la ligne, entre deux traits ; « même
propriété » est souligné}
\marginal{et \uncertain{sont} \uncertain{toutes} de la \uncertain{forme}
$(p^{\frac{1}{2}})$ \ill{} $\xi_i$, où $\nu$ un \uncertain{entier},
$\xi_i$ un \emph{\ill{}} \ill{}}
\note{note écrite verticalement en marge gauche, entourée d'un trait qui
la rattache à « $L(0) = 1$ » ; l'exposant de $p^{\frac{1}{2}}$ est
couvert d'une tache (sans doute le $\nu$ qui suit)}
\[
  L(t) = L_{\rho}(t) \cdots L_{\rho+n}(t)\,
  L'_{\rho+n-1}(pt)\, L'_{\rho+n-2}(p^{2}t) \cdots L'_{\rho}(p^{n}t)
\]
\note{l'indice $\rho+n$ surcharge un autre indice ; une accolade souligne
$L_{\rho}(t) \cdots L_{\rho+n}(t)$ et, au-dessous, un signe barré
(peut-être « $\neq$ ») ne se lit pas ; les arguments $p^{2}t$ et $p^{n}t$
se lisent mal (l'exposant pourrait être un $i$)}
où les valeurs propres inverses des $L_i$ \struck{et des} \uncertain{sont}
des entiers alg.\ de val.\ absolue $p^{\frac{i}{2}}$,
et de même pour $L'_i$ (donc $L'_{\rho+n-i}(p^{i}t)$ a des valeurs propres
inverses qui sont des entiers alg.\ de val.\ absolue
\[
  p^{\frac{\rho+n-i}{2}}\, p^{i} = p^{\frac{\rho+n+i}{2}} .
\]

\page{66}
Ceci posé, on \struck{\ill{}} se permet de « corriger » par des $\lambda$
soumis à la forme
\[
  \begin{array}{l}
  \lambda_{\rho}(t) \cdots \lambda_{\rho+n-1}(t)\,
  \lambda_{\rho+n-2}(pt)\, \lambda_{\rho+n-1}(pt)\, \lambda_{\rho+n-2}(p^{2}t) \\
  \qquad\qquad \cdots \lambda_{\rho}(p^{n}t)
  \end{array}
\]
\note{la suite des indices est transcrite telle qu'elle est écrite ;
l'indice $\rho+n-1$ du deuxième facteur surcharge une autre écriture}
donc de poids entre $\rho$ et $\rho+n$ comme $L$, mais de \emph{niveau}
$\leq \rho+n-1$ (et non seulement $\leq \rho+n$ comme $L$).

\struck{Si \ill{} \ill{} $L_{\rho+n}$}
Nous supposons de plus que $A(t)$ lui-même est de \uncertain{même} type.
\[
  \begin{array}{l}
  L\Bigl(\frac{1}{p^{n+\rho}t}\Bigr)
  = L'_{\rho}\Bigl(\frac{1}{p^{\rho}t}\Bigr) L'_{\rho+1}\Bigl(\frac{1}{p^{\rho+1}t}\Bigr)
    \cdots L'_{\rho+n-1}\Bigl(\frac{1}{p^{n+\rho-1}t}\Bigr) \\[6pt]
  \qquad\qquad L_{\rho+n}\Bigl(\frac{1}{p^{\rho+n}t}\Bigr)
    L_{\rho+n-1}\Bigl(\frac{1}{p^{n+\rho}t}\Bigr) \cdots L_{\rho}\Bigl(\frac{1}{p^{n+\rho}t}\Bigr)
  \end{array}
\]
De plus, on suppose que chaque $L_i$ et $L'_i$ a une équation
fonctionnelle du type correspondant
\note{un trait courbe relie ce membre de gauche aux lignes qui suivent}
\[
  \begin{array}{l}
  = \delta'_{\rho}\, t^{\chi'_{\rho}} L'_{\rho}(t) \cdots
    \delta'_{\rho+n-1}\, t^{\chi'_{\rho+n-1}} L'_{\rho+n-1}(t) \\[4pt]
  \qquad \delta_{\rho+n}\, t^{\chi_{\rho}} L_{\rho+n}(t)\,
    \delta_{\rho+n-1}\, t^{\chi_{\rho+n-1}} L_{\rho+n-1}(pt) \cdots
    \delta_{\rho}\, t^{\chi_{\rho}} L_{\rho}(p^{n}t) \\[4pt]
  = \delta\, t^{\chi} L'_{\rho}(t) \cdots L'_{\rho+n-1}(t)\,
    L_{\rho+n}(t)\, L_{\rho+n-1}(pt) \cdots L_{\rho}(p^{n}t)
  \end{array}
\]
\note{devant $\delta_{\rho+n}$, un $L$ est biffé ; l'exposant $\chi_{\rho}$
de ce facteur est écrit ainsi (on attendrait $\chi_{\rho+n}$) ; le premier
$\delta'_{\rho}$ surcharge une autre lettre ; dans la dernière ligne,
quelques tirets sont tracés au-dessus de $(t)$, et l'argument $(pt)$ est
surchargé}

\page{68}
\struck{$A$ \ill{}}
\[
  \frac{L(\frac{1}{qt})}{L(t)}
  = \delta\, t^{\chi}
  \Bigl[\frac{L'_{\rho}(t)}{L'_{\rho}(p^{n}t)} \cdots
    \frac{L'_{\rho+n-1}(t)}{L'_{\rho+n-1}(pt)}\Bigr]
  \Bigl[\frac{L_{\rho+n-1}(pt)}{L_{\rho+n-1}(t)} \cdots
    \frac{L_{\rho}(p^{n}t)}{L_{\rho}(t)}\Bigr]
\]
\note{avant les derniers points de suspension, un $L$ est biffé ; l'indice
du dénominateur $L_{\rho+n-1}(t)$ surcharge une autre écriture}
\[
  A(t) = \delta\, t^{\chi}
  \prod_i \frac{L'_i(t)}{L'_i(p^{\rho+n-i}t)}
  \prod_i \frac{L_i(p^{\rho+n-i}t)}{L_i(t)}
\]
\note{dans le numérateur du second produit, l'argument surcharge un $p$
biffé}

$n = 1$
\[
  \underbrace{L_{\rho}(t)}\, L_{\rho+1}(t)\, \underbrace{L'_{\rho}(pt)}
\]
\[
  L_{\rho}(t)\, L'_{\rho}(pt)
\]
\note{le reste de la page est blanc}

\end{document}
